\documentclass[../../main.tex]{subfiles}
\begin{document}

{\bf [解]}

与えられた領域
\begin{align}
  \begin{cases}
    x+y+z \leqq n \\
    -x+y-z \leqq n \\
    x-y-z \leqq n \\
    -x-y+z \leqq n
  \end{cases}
\end{align}
について考える．

$z = k \in \mathbb{Z}$ と固定すると，
\begin{align}
  \begin{cases}
    x + y \le n - k \\
    -x + y \le n + k \\
    x - y \le n + k \\
    -x - y \le n - k
  \end{cases}
  \iff
  \begin{cases}
    -(n-k) \le x + y \le n - k \\
    -(n+k) \le x - y \le n + k
  \end{cases} \label{eq:1}
\end{align}
である．

これをみたす $(x, y)$ の存在条件から，
\begin{align}
  \begin{cases}
    -(n-k) \le n - k \\
    -(n+k) \le n + k
  \end{cases}
  \therefore
  -n \le k \le n \label{eq:2}
\end{align}
である．
このもとで\cref{eq:1}を図示すると下図斜線部である(境界含む)．

\begin{figure}[htb]
  \centering
  \begin{tikzpicture}[scale=1.55, >=stealth, every node/.style={font=\small}]

    % --- 1. 定数・パラメータ (n = 3, k = 1 の例) ---
    \def\nVal{3.0}
    \def\kVal{1.0}

    \coordinate (O) at (0, 0);

    % 外枠正方形 AGEC の頂点 (手書きスケッチの頂点ラベル配置)
    % A(左上), C(右上), E(右下), G(左下)
    \coordinate (A_sq) at (-\nVal, \nVal);
    \coordinate (C_sq) at (\nVal, \nVal);
    \coordinate (E_sq) at (\nVal, -\nVal);
    \coordinate (G_sq) at (-\nVal, -\nVal);

    % 4つの切断頂点 B, D, F, H
    \coordinate (B) at (-\kVal, \nVal);
    \coordinate (D) at (\nVal, -\kVal);
    \coordinate (F) at (\kVal, -\nVal);
    \coordinate (H) at (-\nVal, \kVal);

    % --- 2. 内部領域のハッチング塗りつぶし ---
    \fill[pattern=north east lines, pattern color=blue!65, opacity=0.75]
    (B) -- (D) -- (F) -- (H) -- cycle;
    \draw[very thick, blue!85!black] (B) -- (D) -- (F) -- (H) -- cycle;

    % --- 3. 外枠正方形の描画 ---
    \draw[very thick, gray!80!black] (A_sq) -- (C_sq) -- (E_sq) -- (G_sq) -- cycle;

    % --- 4. 座標軸 ---
    \draw[->, thick, gray!80!black] (-4.2, 0) -- (4.4, 0) node[right] {$x$};
    \draw[->, thick, gray!80!black] (0, -4.2) -- (0, 4.4) node[above] {$y$};
    \node[below left=2pt] at (O) {$O$};

    % --- 5. 各辺の境界直線の方程式ラベル ---
    % (i) 辺 HB: y = x + n + k (傾き 1)
    \node[purple!90!black, font=\scriptsize\bfseries, rotate=45, above left]
    at ($ (H)!0.5!(B) $) {$y = x + n + k$};

    % (ii) 辺 BD: x + y = n - k <=> y = -x + n - k (傾き -1)
    \node[teal!90!black, font=\scriptsize\bfseries, rotate=-45, above right]
    at ($ (B)!0.5!(D) $) {$x + y = n - k$};

    % (iii) 辺 DF: y = x - (n + k) (傾き 1)
    \node[purple!90!black, font=\scriptsize\bfseries, rotate=45, below right]
    at ($ (D)!0.5!(F) $) {$y = x - (n + k)$};

    % (iv) 辺 FH: x + y = -(n - k) <=> y = -x - (n - k) (傾き -1)
    \node[teal!90!black, font=\scriptsize\bfseries, rotate=-45, below left]
    at ($ (F)!0.5!(H) $) {$x + y = -(n - k)$};

    % --- 6. 頂点からの投影破線と内部座標 (-k, k) ---
    \draw[dashed, gray!60] (B) -- (-\kVal, 0);
    \draw[dashed, gray!60] (H) -- (0, \kVal);
    \draw[dashed, gray!60] (D) -- (0, -\kVal);
    \draw[dashed, gray!60] (F) -- (\kVal, 0);

    % 内部の -k, k ラベル
    \node[right=2pt, font=\scriptsize, teal!90!black] at (0, \kVal) {$k$};
    \node[below=2pt, font=\scriptsize, teal!90!black] at (-\kVal, 0) {$-k$};
    \node[right=2pt, font=\scriptsize, teal!90!black] at (0, -\kVal) {$-k$};
    \node[above=2pt, font=\scriptsize, teal!90!black] at (\kVal, 0) {$k$};

    % --- 7. 頂点プロットとラベル ---
    % 4隅の正方形頂点 (A, C, E, G)
    \node[above left=2pt, font=\footnotesize\bfseries]  at (A_sq) {$A$};
    \node[above right=2pt, font=\footnotesize\bfseries] at (C_sq) {$C$};
    \node[below right=2pt, font=\footnotesize\bfseries] at (E_sq) {$E$};
    \node[below left=2pt, font=\footnotesize\bfseries]  at (G_sq) {$G$};

    % 切断境界の頂点 (B, D, F, H)
    \fill[black] (B) circle (1.8pt) node[above=2pt, font=\footnotesize\bfseries] {$B$};
    \fill[black] (D) circle (1.8pt) node[right=2pt, font=\footnotesize\bfseries] {$D$};
    \fill[black] (F) circle (1.8pt) node[below=2pt, font=\footnotesize\bfseries] {$F$};
    \fill[black] (H) circle (1.8pt) node[left=2pt, font=\footnotesize\bfseries]  {$H$};

    % 正方形の辺長ラベル n, -n
    \node[above=2pt, font=\scriptsize] at (0, \nVal) {$n$};
    \node[below=2pt, font=\scriptsize] at (0, -\nVal) {$-n$};
    \node[below=2pt, font=\scriptsize] at (\nVal, 0) {$n$};
    \node[below=2pt, font=\scriptsize] at (-\nVal, 0) {$-n$};

    % 4隅の直角二等辺三角形の各辺の寸法 (n-k と n+k)
    % 上側
    \node[teal!90!black, font=\scriptsize\bfseries] at ($ (A_sq)!0.5!(B) + (0, 0.25) $) {$n-k$};
    \node[teal!90!black, font=\scriptsize\bfseries] at ($ (A_sq)!0.5!(H) + (-0.35, 0) $) {$n-k$};
    \node[orange!95!black, font=\scriptsize\bfseries] at ($ (C_sq)!0.5!(B) + (0, 0.25) $) {$n+k$};
    \node[orange!95!black, font=\scriptsize\bfseries] at ($ (C_sq)!0.5!(D) + (0.35, 0) $) {$n+k$};

    % 下側
    \node[teal!90!black, font=\scriptsize\bfseries] at ($ (E_sq)!0.5!(D) + (0.35, 0) $) {$n-k$};
    \node[teal!90!black, font=\scriptsize\bfseries] at ($ (E_sq)!0.5!(F) + (0, -0.25) $) {$n-k$};
    \node[orange!95!black, font=\scriptsize\bfseries] at ($ (G_sq)!0.5!(F) + (0, -0.25) $) {$n+k$};
    \node[orange!95!black, font=\scriptsize\bfseries] at ($ (G_sq)!0.5!(H) + (-0.35, 0) $) {$n+k$};

  \end{tikzpicture}
  \caption{平面 $z=k$ による切断面における条件を満たす領域}
  \label{fig:square_four_triangles_decomposition_labeled_with_lines}
\end{figure}

この内部にある格子点の数を $S(k)$ とおく．
ここで1辺 $2n$ の正方形から，三角形4つを取り去ることで $S(k)$ をもとめる．

まず一辺$2n$の正方形の周及び内部の格子点の数は
\begin{align}
  (2n + 1)^2 \label{eq:3}
\end{align}
である．

\begin{figure}[htb]
  \centering
  \begin{tikzpicture}[scale=1.6, >=stealth, every node/.style={font=\small}]

    % --- 1. 定数・幾何パラメータ ---
    % 一辺 p (例: p = 3.6)
    \def\pVal{3.6}
    \def\kVal{1.4}
    \pgfmathsetmacro{\yOnLine}{\pVal - \kVal} % 2.2

    \coordinate (O)     at (0, 0);
    \coordinate (Pt_pY) at (0, \pVal);
    \coordinate (Pt_pX) at (\pVal, 0);
    \coordinate (Pt_kX) at (\kVal, 0);
    \coordinate (Pt_kLine) at (\kVal, \yOnLine);

    % --- 2. 座標軸 ---
    \draw[->, thick, gray!80!black] (-0.4, 0) -- (4.4, 0) node[right] {$x$};
    \draw[->, thick, gray!80!black] (0, -0.4) -- (0, 4.4) node[above] {$y$};
    \node[below left=2pt] at (O) {$O$};

    % --- 3. 直角二等辺三角形の境界線 ---
    % 斜辺 x + y = p
    \draw[very thick, red!85!black] (Pt_pY) -- (Pt_pX);
    \node[red!90!black, font=\footnotesize\bfseries, above right]
    at ($ (Pt_pY)!0.45!(Pt_pX) $) {$x + y = p$};

    % 軸上の境界
    \draw[thick, gray!70] (0, 0) -- (0, \pVal);
    \draw[thick, gray!70] (0, 0) -- (\pVal, 0);

    % --- 4. 斜辺上の除外点（白抜き丸）---
    % 斜辺上の格子点は除外されるため白抜き丸で明示
    \draw[thick, red!85!black, fill=white] (Pt_pY) circle (2.2pt);
    \draw[thick, red!85!black, fill=white] (Pt_pX) circle (2.2pt);
    \draw[thick, red!85!black, fill=white] (Pt_kLine) circle (2.2pt);

    % その他の斜辺上のサンプル点（白抜き丸）
    \draw[thick, red!85!black, fill=white] (0.6, {\pVal - 0.6}) circle (1.8pt);
    \draw[thick, red!85!black, fill=white] (2.8, {\pVal - 2.8}) circle (1.8pt);

    % --- 5. x = 0 における格子点 (p 個) ---
    % 含まれる点: (0, 0), (0, 1), ..., (0, p-1)
    \fill[black] (0, 0) circle (1.8pt);
    \fill[black] (0, 1.0) circle (1.8pt);
    \fill[black] (0, {\pVal - 0.5}) circle (1.8pt);
    \node[left=3pt, font=\scriptsize] at (0, {\pVal - 0.5}) {$p-1$};
    \node[left=3pt, font=\scriptsize] at (Pt_pY) {$p$};

    % --- 6. 任意の位置 x = k における格子点数のカウント ---
    % 垂直破線
    \draw[dashed, teal!85!black, thick] (Pt_kX) -- (Pt_kLine);
    \fill[black] (Pt_kX) circle (1.8pt);
    \node[below=2pt, font=\scriptsize] at (Pt_kX) {$k$};

    % 含まれる格子点の例 (黒丸)
    \fill[teal!85!black] (\kVal, 0.6) circle (1.8pt);
    \fill[teal!85!black] (\kVal, 1.2) circle (1.8pt);
    \fill[teal!85!black] (\kVal, {\yOnLine - 0.5}) circle (1.8pt);

    % 手書きの波括弧（引き出し線）
    \draw[->, thick, teal!90!black]
    ($ (\kVal, 0.9) + (-0.3, 0) $) to[out=180, in=90]
    (0.6, -0.6) node[below=1pt, font=\footnotesize\bfseries, align=left]
    {$p - k$ 個の格子点\\ \scriptsize ($y = 0, 1, \dots, p-k-1$)};

    % 目盛り p
    \node[below=2pt, font=\scriptsize] at (Pt_pX) {$p$};

  \end{tikzpicture}
  \caption{直角二等辺三角形の周および内部（斜辺を除く）にある格子点の数え方}
  \label{fig:triangle_lattice_counting}
\end{figure}

次に三角形の周および内部（斜辺上を除く）の格子点の数を数える．
上図の $\triangle ABH, \triangle DEF$ は一辺が$n-k$の，$\triangle BCD, \triangle GFH$は一辺が$n+k$の直角二等辺三角形である．一般に一辺$p \in \mathbb{N}$ の直角2等辺三角形の周及び内部(斜辺をのぞく)の格子点の数は
\begin{align}
  1 + 2 + \dots + p &= \dfrac{1}{2} p (p + 1)
\end{align}
である．よって，$\triangle ABH, DEF, BCD, FGH$ の周及び内部(斜辺除く)にある格子点はあわせて
\begin{align}
  2\cdot\dfrac{(n-k)(n+1-k)}{2} + 2\cdot\dfrac{(n+k)(n+k+1)}{2}
  = (n+k)(n+k+1) + (n-k)(n+1-k) \label{eq:4}
\end{align}
である．\cref{eq:3,eq:4}から$S(k)$は
\begin{align}
  S(k) &= (2n+1)^2 - (n+k)(n+1+k) - (n-k)(n+1-k) \label{eq:5}
\end{align}
である．よって求める格子点の合計は\cref{eq:2}より
\begin{align}
  f(n)
  &= \sum_{k=-n}^{n} S(k) \\
  &= S(0) + 2 \sum_{k=1}^n S(k) \\
  &= (2n^2+2n+1) + 2\sum_{k=1}^{n}(2n+1)^2 - 2\sum_{k=1}^{n}(n+k)(n+1+k) - 2\sum_{k=1}^{n}(n-k)(n+1-k) \label{eq:6}
\end{align}
である．各項計算すると
\begin{align}
  (\text{第二項})
  &= 2\sum_{k=1}^{n}(2n+1)^2 \\
  &= 2n(2n+1)^2 \label{eq:7}
\end{align}
\begin{align}
  (\text{第三項})
  &= \sum_{k=1}^n (k+n)(k+1+n) \\
  &= \dfrac{1}{3} \sum_{k=1}^n \left\{ (k+n)(k+1+n)(k+n+2) - (k+n-1)(k+n)(k+n+1) \right\} \\
  &= \dfrac{1}{3} \left\{ (2n)(2n+1)(2n+2) - n(n+1)(n+2) \right\} \\
  &= \dfrac{1}{3} n(n+1)(7n + 2) \label{eq:8}
\end{align}
\begin{align}
  (\text{第四項})
  &=  \sum_{k=1}^n (k-n)(k-n-1) \\
  &= \dfrac{1}{3} \sum_{k=1}^n \left\{ (k+1-n)(k-n)(k-1-n) - (k-n)(k-1-n)(k-2-n) \right\} \\
  &= \dfrac{1}{3} \left\{ -(-1)(-n)(-1-n) \right\} \\
  &= \dfrac{1}{3} n(n+1)(n-1) \label{eq:9}
\end{align}
である．

\cref{eq:7,eq:8,eq:9}を\cref{eq:6}に代入して
\begin{align}
  f(n)
  &= (2n^2+2n+1) +  2n(2n+1)^2 - \dfrac{2}{3} n(n+1)(7n+2) - \dfrac{2}{3} n(n+1)(n-1)
\end{align}
である．両辺を$n^3$で割って
\begin{align}
  \dfrac{f(n)}{n^3}
  &= \dfrac{2n^2+2n+1}{n^3} + 2 \left[ \left(2 + \dfrac{1}{n}\right)^2 - \dfrac{1}{3}\left(1 + \dfrac{1}{n}\right)\left(7 + \dfrac{2}{n}\right) - \dfrac{1}{3}\left(1 + \dfrac{1}{n}\right)\left(1 - \dfrac{1}{n}\right) \right] \\
  &\xrightarrow{n \to \infty} 0 + 2 \left[ 4 - \dfrac{7}{3} - \dfrac{1}{3} \right] \\
  &= \dfrac{8}{3}
\end{align}
が求める極限値である．

\newpage

\end{document}
